\chapter{SIMD：从标量循环到向量化}

SIMD 的全称是 Single Instruction Multiple Data，即一条指令处理多个数据 lane。对 AI 计算而言，SIMD 不是装饰，而是单核吞吐的核心来源。一个标量循环每次处理一个 float，向量化后一次可以处理四个、八个、十六个甚至更多元素。问题在于，数学上能并行，不代表代码一定能被编译器或手写 intrinsic 正确高效地向量化。

\section{向量化需要独立性}

最容易向量化的是逐元素循环。例如 $C_i=A_i+B_i$，每个元素的计算互不依赖，地址连续，循环次数清楚。编译器可以把多个元素合成一个向量加载、一个向量加法和一个向量存储。若循环中存在跨迭代依赖，例如前缀和，后一个元素依赖前一个元素的结果，向量化就困难得多。

\begin{lstlisting}[language=C++,caption={容易向量化的逐元素循环}]
#include <cstdint>
#include <span>

void add_arrays(std::span<const float> a,
                std::span<const float> b,
                std::span<float> c) {
    const std::int64_t n = static_cast<std::int64_t>(c.size());
    for (std::int64_t i = 0; i < n; ++i) {
        c[static_cast<std::size_t>(i)] =
            a[static_cast<std::size_t>(i)] + b[static_cast<std::size_t>(i)];
    }
}
\end{lstlisting}

这段代码看起来普通，但它具备向量化的基本条件：连续内存、固定步长、没有循环携带依赖。若把输出写到输入重叠的区域，或者通过不透明指针访问，编译器就可能因为别名风险保守。工业代码会通过接口约束、对齐信息、restrict 风格提示或库内部封装减少这种不确定性。

\section{归约为什么更复杂}

点积也是 AI 计算常见循环，但它比逐元素加法复杂，因为所有迭代都更新同一个累加器。向量化点积通常不是让多个 lane 同时写一个标量，而是维护多个部分和，最后再做水平归约。这样做可以减少依赖链长度，提高指令级并行。

\begin{lstlisting}[language=C++,caption={教学版点积：多个累加器降低依赖链压力}]
#include <cstdint>
#include <span>

float dot_unrolled(std::span<const float> a,
                   std::span<const float> b) {
    const std::int64_t n = static_cast<std::int64_t>(a.size());
    float acc0 = 0.0F;
    float acc1 = 0.0F;
    float acc2 = 0.0F;
    float acc3 = 0.0F;

    std::int64_t i = 0;
    for (; i + 3 < n; i += 4) {
        acc0 += a[static_cast<std::size_t>(i)] *
                b[static_cast<std::size_t>(i)];
        acc1 += a[static_cast<std::size_t>(i + 1)] *
                b[static_cast<std::size_t>(i + 1)];
        acc2 += a[static_cast<std::size_t>(i + 2)] *
                b[static_cast<std::size_t>(i + 2)];
        acc3 += a[static_cast<std::size_t>(i + 3)] *
                b[static_cast<std::size_t>(i + 3)];
    }

    float acc = (acc0 + acc1) + (acc2 + acc3);
    for (; i < n; ++i) {
        acc += a[static_cast<std::size_t>(i)] *
               b[static_cast<std::size_t>(i)];
    }
    return acc;
}
\end{lstlisting}

这不是最终的 SIMD intrinsic 版本，但它揭示了一个重要思想：归约的瓶颈不只是操作数量，还有依赖链。现代 CPU 可以同时执行多条独立浮点指令，但如果每条指令都依赖上一条累加结果，吞吐就会被延迟限制。多个累加器让硬件有更多独立工作可做。

\section{对齐、尾部和掩码}

SIMD 循环通常要处理三个细节。第一是对齐。对齐访问更容易让硬件高效工作，虽然现代 CPU 对非对齐访问的支持已经比过去好很多。第二是尾部。数组长度不一定刚好是向量宽度的整数倍，尾部元素要用标量循环、掩码加载或特殊路径处理。第三是异常和边界。若向量加载越过合法内存，即使多出来的 lane 不参与结果，也可能触发错误；因此边界处理必须明确。

AVX-512 提供更强的掩码机制，能让尾部处理更自然；AVX2 常见做法是主循环处理完整向量，尾部用标量补齐。不同平台有不同权衡。本册不会把所有 intrinsic 当成必背 API，而会先建立向量语义：一组 lane 同步执行，同一条指令处理多个元素，水平归约和跨 lane shuffle 是额外成本。

\section{编译器自动向量化与手写内核}

能自动向量化时，应先让编译器做。编译器知道目标架构、寄存器分配和调度细节，也能随编译选项适配不同平台。但矩阵乘、卷积和量化 GEMM 这类核心内核往往需要手写或生成，因为它们涉及寄存器块、数据 packing、预取、指令选择和微架构调度。手写内核不是为了炫技，而是因为通用编译器很难从普通三重循环推导出工业库需要的全部结构。

\begin{warningbox}
不要把 SIMD 等同于“把循环步长改成向量宽度”。真正的问题包括别名、对齐、尾部、依赖链、数据布局、寄存器压力和内存带宽。向量化增加了单核计算能力，也可能把瓶颈更快推到内存系统。
\end{warningbox}

\begin{labbox}
实验建议：使用编译器的向量化报告，例如 GCC 的 \cmd{-fopt-info-vec} 或 Clang 的 \cmd{-Rpass=loop-vectorize}，观察逐元素加法、点积、前缀和三个循环分别能否自动向量化。把报告和汇编对照起来，确认是否出现向量加载和向量浮点指令。
\end{labbox}
